\section{Virtual difference mode after handoff}
\label{sec:virtualmode}

The local encoder does not assume a bosonic mode already entangled with $S$. That bosonic channel input emerges only after the local preparation has completed.

\begin{lemma}[Virtual difference-mode reduction]
After the handoff, suppose the branch-conditioned state of all bosonic degrees of freedom in the vacuum/coherent linear network is, after removing branch-common displacements,
\begin{equation}
|\Psi_s\rangle=|s\bm{\alpha}\rangle,
\qquad s=\pm1.
\end{equation}
Let $A=\|\bm{\alpha}\|$ and define
\begin{equation}
\boxed{
d=\frac{1}{A}\sum_j\alpha_j^*b_j.}
\label{eq:virtualmode}
\end{equation}
Then a branch-independent passive mode rotation maps
\begin{equation}
|s\bm{\alpha}\rangle
\longmapsto
|sA\rangle_d\otimes|0\rangle_\perp.
\label{eq:modecompression}
\end{equation}
For any selected receiver memory mode with coherent branch amplitude $\alpha_B$, the reduced map from $d$ to that memory is a pure-loss projection of transmissivity
\begin{equation}
\boxed{
\eta_B=\frac{|\alpha_B|^2}{A^2}
=\frac{N_{\Delta,B}}{N_{\Delta,{\rm all}}}.}
\label{eq:etaB}
\end{equation}
\end{lemma}

The proof is a unitary rotation in one-particle mode space that aligns the displacement vector $\bm{\alpha}$ with one basis vector. Because coherent states transform covariantly under passive mode rotations, all branch dependence is compressed into $d$ and the orthogonal modes are vacuum. This lemma is the formal bridge between the physical local encoder and the standard binary-coherent bosonic-channel problem.

A branch-common controller or background gravitational coherent field causes only a common displacement, $|\gamma_C\pm\alpha\rangle$, and is removed by the same branch-independent displacement. It does not change $N_\Delta$, NPT, or entanglement-breaking status. Controller fluctuations above the coherent-state vacuum level are different: they are genuine added noise and must be propagated through the channel.

\section{Source gravitational branching and free-space propagation}
\label{sec:sourceprop}

Let the source mode couple to independent narrowband Markov output ports $j$ with rates $\kappa_j$. For a mirrored coherent trajectory $\alpha_s(t)=s\alpha(t)$, the branch-difference output amplitude is
\begin{equation}
\Delta b_j^{\rm out}(t)=2\sqrt{\kappa_j}\,\alpha(t).
\end{equation}
Hence
\begin{equation}
N_{\Delta,j}=4\kappa_j\int\dd t\,|\alpha(t)|^2.
\end{equation}
The fraction of the complete branch-distance norm emitted into the gravitational port is therefore
\begin{equation}
\boxed{
\beta_{g,A}
\equiv\frac{N_{\Delta,g}}{N_{\Delta,{\rm tot}}}
=\frac{\kappa_{g,A}}{\kappa_A},}
\qquad
\kappa_A=\sum_j\kappa_j.
\label{eq:betaA}
\end{equation}
This equality is a narrowband amplitude-damping result. If the port couplings vary appreciably across the source waveform bandwidth, the correct object is the waveform-weighted ratio
\begin{equation}
\beta_{g,A}[\alpha]
=
\frac{\int\dd\Omega\,\kappa_g(\omega+\Omega)|\widetilde\alpha(\Omega)|^2}
{\int\dd\Omega\,\kappa_{\rm tot}(\omega+\Omega)|\widetilde\alpha(\Omega)|^2}.
\label{eq:freqbeta}
\end{equation}

Equation~\eqref{eq:betaA} is not a generic coherence parameter. Pure phase diffusion can destroy the source-reference entanglement without changing the source energy or any branching ratio. Non-Gaussian source dephasing must therefore be represented by a separate source channel rather than hidden inside $\beta_{g,A}$.

\subsection{Wave-zone mode capture}

For aligned compact plus quadrupoles in the wave zone, the retarded source-to-receiver normalization gives the useful receiver gravitational-port rate
\begin{equation}
\kappa_\Delta(R)=\eta_{\rm store}(R)\kappa_{g,B},
\end{equation}
with
\begin{equation}
\boxed{
\eta_{\rm store}(R)=\frac{25\Ocal}{16(kR)^2},
\qquad k=\omega/c,}
\label{eq:etastore}
\end{equation}
where $0\le\Ocal\le1$ collects residual temporal, tensor, polarization, and orientation mode overlap.

The coefficient has a reciprocal-antenna interpretation. The on-axis directivity of the aligned plus source is $G_A=5/2$, while the critically coupled $l=2$ receiver effective area $A_e=5\pi/(2k^2)$ corresponds to reciprocal gain $G_B=5/2$. Thus
\begin{equation}
\boxed{
\eta_{\rm store}
=\Ocal G_AG_B\left(\frac{\lambda}{4\pi R}\right)^2
=\frac{25\Ocal}{16(kR)^2}.}
\label{eq:friis}
\end{equation}
This is ordinary reciprocal far-field link physics specialized to the quadrupole channel, not a new antenna theorem. Critical-coupling and temporal coupled-mode normalization are standard \cite{RuanFan2012}.

\section{Receiver memory and the four-factor coherent link}
\label{sec:memory}

Let the receiver memory mode $b$ have total linewidth $\kappa_B$ and intrinsic gravitational linewidth $\kappa_{g,B}$. Define
\begin{equation}
\boxed{\beta_{g,B}=\frac{\kappa_{g,B}}{\kappa_B}.}
\label{eq:betaB}
\end{equation}
After retarded arrival, let $f(t)$ be the normalized temporal envelope of the selected incoming source mode,
\begin{equation}
\int_0^\infty\dd t\,|f(t)|^2=1.
\end{equation}
Here $f$ is the complete source gravitational waveform: it may contain encoder precursor radiation emitted before the physical handoff as well as later source radiation. At any chosen receiver observation time, the convolution below includes whatever portion of that complete waveform has arrived by then.

The linear receiver equation contains the selected gravitational input with rate $\kappa_\Delta=\eta_{\rm store}\kappa_{g,B}$. The coherent loading from the selected mode is therefore
\begin{equation}
\tau_{f,B}(t)
=\kappa_\Delta
\left|
\int_0^t\dd s\,
\ee^{-\kappa_B(t-s)/2}f(s)
\right|^2.
\end{equation}
Define the dimensionless temporal loading factor
\begin{equation}
\boxed{
\Tcal_f(t)
=\kappa_B
\left|
\int_0^t\dd s\,
\ee^{-\kappa_B(t-s)/2}f(s)
\right|^2,}
\qquad
0\le\Tcal_f(t)\le1.
\label{eq:Tf}
\end{equation}
The upper bound follows from Cauchy--Schwarz and is saturated at a chosen target time by the time reverse of the receiver impulse response.

\begin{proposition}[Post-handoff gravitational quantum-link budget]
For the vacuum/coherent, narrowband linear amplitude-damping model, once the full source/controller system has passed its physical handoff, the completed logical code defines the fixed virtual difference mode of Sec.~\ref{sec:virtualmode}. At any subsequent receiver observation time, with the complete waveform $f$ including encoder precursor radiation, the receiver branch-distance fraction relative to that completed code is
\begin{equation}
\boxed{
\tau_c(t)=
\beta_{g,A}\,
\eta_{\rm store}(R)\,
\beta_{g,B}\,
\Tcal_f(t).}
\label{eq:master}
\end{equation}
Defining
\begin{equation}
\boxed{\eta_Q^{\rm link}=\beta_{g,A}\eta_{\rm store}\beta_{g,B},}
\label{eq:etalink}
\end{equation}
one has $\tau_c(t)=\eta_Q^{\rm link}\Tcal_f(t)$ and $\tau_c\le\eta_Q^{\rm link}$.
\end{proposition}

The factors in Eq.~\eqref{eq:master} are not duplicate normalizations. The same product follows from an independent cascaded bosonic network: a source beamsplitter of transmissivity $\beta_{g,A}$, a propagation beamsplitter of transmissivity $\eta_{\rm store}$, and a receiver loading channel whose coherent fraction is $\beta_{g,B}\Tcal_f$. The product is therefore the canonical serial loss of the post-handoff virtual mode.

If the receiver has already begun responding to precursor radiation before the physical source/controller handoff, the same normalized receiver response coefficient can be computed from the instantaneous coherent branch state, but it should not yet be interpreted as the transmissivity of a fixed logical-to-bosonic input channel. The fixed-input channel description begins only once the controller degrees have been cleared and the completed branch-distance norm is established.

Two useful waveform benchmarks illustrate the order-unity loading factor. A matched passive exponential gives
\begin{equation}
\Tcal_{\rm exp}^{\max}=4\ee^{-2}\simeq0.541341,
\end{equation}
while target-time optimization yields
\begin{equation}
\Tcal_{\rm opt}(t)=1-\ee^{-\kappa_B t}\to1.
\end{equation}
Temporal shaping can remove an order-unity mismatch penalty, but it cannot exceed the source-branching, propagation, or receiver-branching ceiling.
